
We set out to ask whether the limited value of in-season replanning reported
for irrigation survives when the biological outcome responds sharply to
timing error. It does not, and the boundary can be located.

The gain from adaptive replanning rises with the curvature of the yield
response on every farm tested ($r = 0.94$), crossing the five percent
mark at about 1.8 times the curvature measured for hybrid maize.
Below that the static plan absorbs the risk; above it, it cannot. The
mechanism is hedging. A sowing offset is a decision taken once with a tunable
margin, so when the penalty steepens the static optimiser responds by
choosing an offset that stays tolerable across the plausible range of
thermal-time outcomes rather than aiming at the peak. That substitution works
until the penalty becomes steep enough that no single margin is tolerable
everywhere, and it is at that point that acting on new information starts to
pay. An irrigation depth, which must be chosen repeatedly and cannot be
hedged in the same way, sits in a different part of this space, which is why
the irrigation literature saw only the flat regime.

Hybrid maize, at the sharpness its own field data imply, falls below the
boundary: adaptivity is worth 1.91\% ($p = 0.003$), real but modest, and
concentrated in the seasons and on the farms that go worst. For this crop the
Kelly--Becker recommendation transfers. The contribution is not that it
transfers, but that we can now say how far it transfers and what takes an
operation outside it.

\subsection{What single-instance designs hide}

Both of the corrections above came from the same fix. Measured on one farm,
the adaptivity gain was 3.9\% and appeared to saturate below five
percent; measured on ten, it is 1.91\% and saturates nowhere. One
instance produced an effect roughly twice its true size and a qualitative
conclusion that reversed on replication.

We draw attention to this because simulation studies in this area routinely
report results from a single synthetic instance, and because neither error
was visible from inside that instance: the single-farm result was internally
consistent, adequately powered across seasons, and statistically significant.
Replication across instances is not a robustness check to be added if space
allows. It is what determines whether the headline number means anything.

\subsection{A screening rule}

Three questions decide whether an adaptive scheduler is worth building.

\emph{How sharp is the timing penalty?} Below roughly 1.8 times the
curvature of the hybrid maize response, a well-built static plan captures
nearly everything available. This is the dominant factor.

\emph{Can the decision be hedged in advance?} A choice made once with a
tunable margin, such as a sowing offset, can absorb uncertainty before the
season; one that must be made repeatedly cannot. Hedging substitutes for
adapting.

\emph{How well does a static plan already do?} The gain was strongly
anti-correlated with baseline performance ($r = -0.75$). An operation
already close to its synchrony optimum has little for adaptation to recover.

\subsection{What the benchmark choice decides}

Had we compared adaptive replanning against conventional practice rather than
against a robust static plan, we would have reported a substantially larger
figure and claimed a strong result. Most of the apparent benefit of adaptive
scheduling in this problem is the benefit of optimising at all, which a static
solver delivers once, before the season, at a fraction of the operational
cost. We also note that both failure modes encountered in development ---
indefinite deferral under free reassignment, and double-counting rain risk in
objective and constraints --- caused the adaptive policy to \emph{underperform}
the static one. An adaptive system is better than a fixed one only if these
are handled.
